\documentclass[11pt,a4paper]{article}

\usepackage[margin=1in]{geometry}
\usepackage[T1]{fontenc}
\usepackage[utf8]{inputenc}
\usepackage{lmodern}
\usepackage{xcolor}
\usepackage{hyperref}
\usepackage{longtable}
\usepackage{array}
\usepackage{booktabs}
\usepackage{enumitem}
\usepackage{listings}
\usepackage{titlesec}

\definecolor{miganavy}{HTML}{16324D}
\definecolor{migablue}{HTML}{2F6F97}
\definecolor{migasoft}{HTML}{EAF3FA}
\definecolor{migatext}{HTML}{31465A}

\hypersetup{
    colorlinks=true,
    linkcolor=miganavy,
    urlcolor=migablue,
    pdftitle={Hardware Controller User Manual},
    pdfauthor={OpenAI Codex},
}

\titleformat{\section}{\Large\bfseries\color{miganavy}}{\thesection}{0.7em}{}
\titleformat{\subsection}{\large\bfseries\color{miganavy}}{\thesubsection}{0.7em}{}

\lstdefinestyle{migacode}{
    basicstyle=\ttfamily\small,
    backgroundcolor=\color{migasoft},
    frame=single,
    framerule=0pt,
    breaklines=true,
    showstringspaces=false,
    columns=fullflexible,
}

\title{\textbf{Hardware Controller User Manual}\\[0.4em]
\large Integrated Multi-Device EDFA and PSU Web Control}
\author{Prepared for the MIGA laboratory workflow}
\date{Revision date: July 27, 2026}

\begin{document}
\maketitle

\tableofcontents
\newpage

\section{Purpose}
This document describes how to use the \texttt{Hardware-controller} application created to integrate two existing laboratory tools:
\begin{itemize}[leftmargin=1.5em]
    \item the EDFA remote control workflow from \texttt{Edfa-controller/edfa-controller.py},
    \item the programmable power supply scheduler from \texttt{Power-supply-remote/Dispenser\_remote\_control-2channels\_STD.py}.
\end{itemize}

The new application exposes both functions through a single local web service. The service can be opened in a browser on the host computer or from another computer on the same local area network (LAN), provided the host and network permissions allow it.

\section{Project Location and File Layout}
The project lives inside:

\begin{lstlisting}[style=migacode]
/home/yiming/Documents/MIGA-program/Hardware-controller
\end{lstlisting}

The main files are listed in Table~\ref{tab:layout}.

\begin{longtable}{>{\raggedright\arraybackslash}p{0.31\linewidth} >{\raggedright\arraybackslash}p{0.61\linewidth}}
\caption{Important files in the Hardware Controller project.}\label{tab:layout}\\
\toprule
\textbf{Path} & \textbf{Role} \\
\midrule
\endfirsthead
\toprule
\textbf{Path} & \textbf{Role} \\
\midrule
\endhead
\texttt{main.py} & Uvicorn/FastAPI entry point. \\
\texttt{config.py} & Default ports, timing constants, password policy, and file paths. \\
\texttt{app/\_\_init\_\_.py} & Builds the FastAPI application, login page, main page, and WebSocket bridge. \\
\texttt{app/api/routes.py} & Browser-facing HTTP API for login, device management, manual actions, and downloads. \\
\texttt{app/core/controller\_manager.py} & Main coordination layer for devices, scheduling, state updates, and event logging. \\
\texttt{app/core/scheduler.py} & Background scheduler thread for EDFA and PSU timed actions. \\
\texttt{app/core/state\_store.py} & Persistent JSON storage for devices, schedules, and password hash. \\
\texttt{app/core/auth.py} & Session handling and password change logic. \\
\texttt{app/drivers/edfa.py} & TCP command sender for EDFA devices. \\
\texttt{app/drivers/psu.py} & Persistent socket driver for the QPX-like power supply protocol. \\
\texttt{static/index.html} & Main browser control console. \\
\texttt{static/login.html} & Login page. \\
\texttt{static/app.js} & Frontend logic for browser actions and live updates. \\
\texttt{static/styles.css} & Frontend styling. \\
\texttt{data/controller\_state.json} & Persistent runtime configuration, device list, schedules, and password hash. \\
\texttt{docs/hardware\_controller\_manual.tex} & This manual. \\
\bottomrule
\end{longtable}

\section{System Design Summary}
The architecture follows the same general web-service idea used in the existing \texttt{MIGA-controller} project:
\begin{itemize}[leftmargin=1.5em]
    \item a Python backend using \texttt{FastAPI},
    \item browser pages served locally by the same backend,
    \item a JSON state file to persist configuration across restarts,
    \item a WebSocket channel to push event messages to the browser,
    \item a background scheduler thread to handle timed switching.
\end{itemize}

The controller does \textbf{not} depend on an external cloud service. It is designed for local laboratory use.

\section{Hardware Capabilities}
\subsection{EDFA Systems}
The EDFA integration assumes the same control logic as the original EDFA script:
\begin{itemize}[leftmargin=1.5em]
    \item each EDFA system exposes four channels,
    \item the channel names are \texttt{edfa0}, \texttt{edfa1}, \texttt{edfa2}, and \texttt{edfa3},
    \item the default startup powers are:
    \begin{itemize}[leftmargin=1.5em]
        \item \texttt{edfa0 = 3},
        \item \texttt{edfa1 = 2.4},
        \item \texttt{edfa2 = 3},
        \item \texttt{edfa3 = 3}.
    \end{itemize}
    \item the browser allows these power values to be edited per device,
    \item a device-wide schedule can turn all four channels on using the stored powers, and off using the safe shutdown sequence.
\end{itemize}

\subsection{Power Supply Systems}
The PSU integration assumes the same protocol as the original scheduler:
\begin{itemize}[leftmargin=1.5em]
    \item each PSU uses TCP port \texttt{9221} by default,
    \item each PSU provides two independently controlled outputs,
    \item channel status can be queried back from the instrument,
    \item independent mode can be applied from the browser,
    \item each channel can have its own weekly schedule and its own on/off times.
\end{itemize}

\section{Installation}
\subsection{Python Dependencies}
The required Python packages are listed in \texttt{requirements.txt}.

Install them with:

\begin{lstlisting}[style=migacode]
cd /home/yiming/Documents/MIGA-program/Hardware-controller
python3 -m pip install -r requirements.txt
\end{lstlisting}

\subsection{Recommended Launch Command}
The default listening port defined in \texttt{config.py} is \texttt{8050}. A practical launch command is:

\begin{lstlisting}[style=migacode]
cd /home/yiming/Documents/MIGA-program/Hardware-controller
python3 -m uvicorn main:app --host 0.0.0.0 --port 8050
\end{lstlisting}

This makes the service available on all network interfaces of the host computer.

\section{Accessing the Service}
\subsection{Local Browser Access}
From the host computer, open:

\begin{lstlisting}[style=migacode]
http://127.0.0.1:8050
\end{lstlisting}

\subsection{LAN Access}
From another computer on the same LAN, open:

\begin{lstlisting}[style=migacode]
http://<host-ip-address>:8050
\end{lstlisting}

Replace \texttt{<host-ip-address>} with the actual IPv4 address of the host computer, for example \texttt{192.168.0.42}.

\subsection{Firewall Considerations}
If LAN access does not work:
\begin{itemize}[leftmargin=1.5em]
    \item check that the host machine firewall allows inbound TCP connections on port \texttt{8050},
    \item check that the browser machine can reach the host on the same subnet,
    \item check that no other service is already occupying port \texttt{8050}.
\end{itemize}

\section{Authentication}
\subsection{Default Password}
The initial login password is:

\begin{lstlisting}[style=migacode]
miga
\end{lstlisting}

\subsection{Password Storage}
The controller stores a password hash and salt inside:

\begin{lstlisting}[style=migacode]
data/controller_state.json
\end{lstlisting}

The plaintext password is not written to disk.

\subsection{Password Change Workflow}
After login:
\begin{enumerate}[leftmargin=1.5em]
    \item scroll to the \textbf{Security \& Service} section,
    \item enter the current password,
    \item enter a new password,
    \item click \textbf{Update Password}.
\end{enumerate}

The current browser session is refreshed automatically after a successful password change.

\section{Persistent State}
The application automatically saves:
\begin{itemize}[leftmargin=1.5em]
    \item all EDFA device entries,
    \item all PSU device entries,
    \item all schedule settings,
    \item device notes,
    \item the password hash,
    \item recent device runtime state such as last action and last communication error.
\end{itemize}

The persistent file is:

\begin{lstlisting}[style=migacode]
Hardware-controller/data/controller_state.json
\end{lstlisting}

When the service restarts, the file is loaded automatically.

\section{Using the Main Dashboard}
\subsection{Top Bar}
The top bar displays:
\begin{itemize}[leftmargin=1.5em]
    \item WebSocket connection status,
    \item scheduler status,
    \item a button to download this \LaTeX{} manual,
    \item a logout button.
\end{itemize}

\subsection{Sidebar}
The sidebar provides:
\begin{itemize}[leftmargin=1.5em]
    \item a summary of EDFA and PSU device counts,
    \item EDFA batch actions,
    \item the \textbf{Apply EDFA Template} form,
    \item a live activity log.
\end{itemize}

\subsection{Main Sections}
The main page contains four sections:
\begin{enumerate}[leftmargin=1.5em]
    \item \textbf{Control Dashboard},
    \item \textbf{EDFA Systems},
    \item \textbf{Power Supply Systems},
    \item \textbf{Security \& Service}.
\end{enumerate}

\section{EDFA Operation}
\subsection{Adding an EDFA Device}
In the \textbf{EDFA Systems} section:
\begin{enumerate}[leftmargin=1.5em]
    \item fill in the device name,
    \item enter the device IP address,
    \item adjust the port, timeout, and command delay if needed,
    \item click \textbf{Add EDFA Device}.
\end{enumerate}

\subsection{Single-Device Commands}
Each EDFA card offers:
\begin{itemize}[leftmargin=1.5em]
    \item \textbf{Save Config},
    \item \textbf{Probe},
    \item \textbf{Start All},
    \item \textbf{Safe Shutdown},
    \item \textbf{Delete}.
\end{itemize}

\subsection{Per-Channel Commands}
Each channel row provides:
\begin{itemize}[leftmargin=1.5em]
    \item a power field,
    \item an \textbf{ON} button,
    \item an \textbf{OFF} button.
\end{itemize}

When \textbf{ON} is pressed, the current power field is used as the command value.

\subsection{Important Note on EDFA State Readback}
The current EDFA protocol in the inherited script provides command sending but no confirmed state query. Therefore:
\begin{itemize}[leftmargin=1.5em]
    \item the UI shows \textbf{assumed} EDFA channel state,
    \item the assumed state is based on the most recent successful controller action,
    \item it should be interpreted as a control log, not as a hardware readback measurement.
\end{itemize}

\subsection{EDFA Weekly Schedule}
Each EDFA device contains one schedule block:
\begin{itemize}[leftmargin=1.5em]
    \item \textbf{Enable schedule},
    \item \textbf{Auto-ON} time,
    \item \textbf{Auto-OFF} time,
    \item active weekdays.
\end{itemize}

When the schedule is enabled:
\begin{itemize}[leftmargin=1.5em]
    \item the ON event applies the stored power preset to all four channels,
    \item the OFF event sends the safe shutdown sequence to all four channels.
\end{itemize}

\subsection{Applying One Template to Many EDFA Devices}
The sidebar contains a batch template form. It can push:
\begin{itemize}[leftmargin=1.5em]
    \item one set of four channel powers,
    \item one schedule definition,
    \item one weekday selection,
\end{itemize}
to all selected EDFA devices, or to all EDFA devices if none are selected.

\section{PSU Operation}
\subsection{Adding a PSU Device}
In the \textbf{Power Supply Systems} section:
\begin{enumerate}[leftmargin=1.5em]
    \item fill in the name and IP address,
    \item keep the default port \texttt{9221} unless your instrument differs,
    \item set the timeout if necessary,
    \item click \textbf{Add PSU Device}.
\end{enumerate}

\subsection{PSU Device Controls}
Each PSU card provides:
\begin{itemize}[leftmargin=1.5em]
    \item \textbf{Save Config},
    \item \textbf{Probe \& IDN},
    \item \textbf{Refresh State},
    \item \textbf{Set Independent Mode},
    \item \textbf{Disconnect},
    \item \textbf{Delete}.
\end{itemize}

\subsection{PSU Channel Controls}
Each PSU offers:
\begin{itemize}[leftmargin=1.5em]
    \item \textbf{CH1 ON} and \textbf{CH1 OFF},
    \item \textbf{CH2 ON} and \textbf{CH2 OFF}.
\end{itemize}

After a successful command, the controller reads the PSU output state back from the instrument and updates the browser.

\subsection{PSU Weekly Schedules}
Each PSU card contains a schedule block for:
\begin{itemize}[leftmargin=1.5em]
    \item channel 1,
    \item channel 2.
\end{itemize}

Each channel schedule contains:
\begin{itemize}[leftmargin=1.5em]
    \item an enable checkbox,
    \item an ON time,
    \item an OFF time,
    \item its own weekday selection.
\end{itemize}

This means the two outputs can run completely independent schedules.

\section{Scheduler Behavior}
The scheduler runs in a background thread inside the Python service. It checks device schedules once per second.

\subsection{Trigger Logic}
\begin{itemize}[leftmargin=1.5em]
    \item A schedule only runs on the selected weekdays.
    \item A schedule action triggers when the current local time matches the configured \texttt{HH:MM}.
    \item A once-per-day guard prevents the same scheduled action from repeating continuously during the same minute.
\end{itemize}

\subsection{Recommended Practice}
For laboratory safety:
\begin{itemize}[leftmargin=1.5em]
    \item verify the host clock before trusting scheduled events,
    \item use simple, clearly separated schedule times,
    \item confirm network reachability for all remote devices before leaving the system unattended,
    \item prefer safe shutdown times for all remotely powered hardware.
\end{itemize}

\section{Event Log}
The sidebar log records:
\begin{itemize}[leftmargin=1.5em]
    \item startup and shutdown events,
    \item device add/update/remove events,
    \item successful commands,
    \item communication failures,
    \item scheduler-triggered actions,
    \item password change events.
\end{itemize}

The log is intended as an operational trace for the active service session.

\section{Troubleshooting}
\subsection{Cannot Reach the Login Page}
Check:
\begin{itemize}[leftmargin=1.5em]
    \item whether \texttt{uvicorn} is running,
    \item whether the chosen port is correct,
    \item whether another service is already using the same port.
\end{itemize}

\subsection{EDFA Probe Fails}
Check:
\begin{itemize}[leftmargin=1.5em]
    \item EDFA IP address,
    \item EDFA port,
    \item device power and network cable,
    \item firewall rules between the browser host and the EDFA host.
\end{itemize}

\subsection{PSU State Refresh Fails}
Check:
\begin{itemize}[leftmargin=1.5em]
    \item PSU IP address,
    \item port \texttt{9221},
    \item whether the PSU network socket is already occupied by another tool,
    \item whether the instrument is in a valid remote control mode.
\end{itemize}

\subsection{Browser Shows Authentication Required}
This usually means:
\begin{itemize}[leftmargin=1.5em]
    \item the login session expired,
    \item the password was changed in another session,
    \item the browser cookie was cleared.
\end{itemize}

Log in again through \texttt{/login}.

\section{Recommended Operational Workflow}
\begin{enumerate}[leftmargin=1.5em]
    \item Start the FastAPI service on the host computer.
    \item Open the login page locally and confirm access.
    \item Change the default password from \texttt{miga}.
    \item Add all EDFA and PSU devices.
    \item Probe each device once and confirm communication.
    \item Save final power presets and schedule definitions.
    \item Test one manual ON/OFF cycle per device before remote unattended use.
    \item If LAN access is required, test the same page from another machine on the subnet.
\end{enumerate}

\section{Suggested Future Extensions}
The current version is intentionally conservative and aligned with the inherited scripts. Reasonable future additions could include:
\begin{itemize}[leftmargin=1.5em]
    \item role-based multi-user authentication,
    \item TLS reverse-proxy deployment,
    \item confirmed EDFA state readback if the hardware protocol later exposes a query command,
    \item configuration export/import workflows,
    \item explicit audit log persistence beyond the active session.
\end{itemize}

\section{Conclusion}
The \texttt{Hardware-controller} project consolidates legacy EDFA and PSU control into one browser-based service while preserving the practical command logic of the original laboratory scripts. It supports multiple EDFA systems, multiple power supplies, persistent device configuration, independent scheduling, LAN access, and a simple password-protected workflow suitable for day-to-day experimental operations.

\end{document}
